% =============================================================================
%                    WEEK 5: TRANSFORMER MODELS AND LLMs
% =============================================================================
\section{Week 5: Transformer Models and Large Language Models}

% -----------------------------------------------------------------------------
%                              5.1 TOPIC SUMMARY
% -----------------------------------------------------------------------------
\subsection{Topic Summary}

\begin{keybox}[Learning Objectives]
By the end of this week, you will be able to:
\begin{enumerate}
  \item Define \textbf{Large Language Models (LLMs)} and explain how they extend earlier language models (n-gram, RNN) through scale and architecture.
  \item Distinguish the three LLM architectures: \textbf{decoder-only} (GPT, Llama), \textbf{encoder-only} (BERT), and \textbf{encoder-decoder} (T5, Whisper).
  \item Explain \textbf{self-supervised pretraining} with cross-entropy loss and teacher forcing on large corpora.
  \item Describe problems with web-scraped training data: quality, bias, privacy, and copyright.
  \item Evaluate LLMs using \textbf{perplexity} and understand its relationship to probability.
  \item Explain why \textbf{static embeddings} (word2vec) are insufficient and how \textbf{contextual embeddings} address this.
  \item Define the \textbf{self-attention} mechanism: queries, keys, values, and scaled dot-product attention.
  \item Describe \textbf{multi-head attention} and why multiple heads are useful.
  \item Explain the components of a \textbf{transformer block}: attention, feedforward layers, residual connections, layer normalisation.
  \item Understand \textbf{causal masking} for autoregressive generation.
  \item Describe \textbf{positional encodings} and why transformers need them.
  \item Explain the \textbf{language modelling head}: unembedding layer, logits, and softmax.
  \item Discuss \textbf{scaling laws} and \textbf{KV caching} for efficient inference.
\end{enumerate}
\end{keybox}

\separator

\subsubsection*{The Big Picture}

This week introduces the \textbf{transformer architecture}, the foundation of modern LLMs. Unlike RNNs which process sequences step-by-step, transformers use \textbf{self-attention} to process all tokens in parallel, enabling much longer contexts and more efficient training. LLMs are transformers trained on massive corpora with billions of parameters, achieving remarkable capabilities through \textbf{pretraining} followed by instruction tuning and alignment.

\separator

\subsubsection*{What Examiners Like to Ask}

\begin{itemize}
  \item Define an LLM and compare it to n-gram/RNN language models.
  \item Explain the three LLM architectures (decoder, encoder, encoder-decoder) and their use cases.
  \item Describe \textbf{self-supervised training}: cross-entropy loss, teacher forcing.
  \item Define \textbf{perplexity} and explain why lower is better.
  \item Explain why static embeddings fail for polysemous words and how attention solves this.
  \item Write the \textbf{scaled dot-product attention} formula and explain each component.
  \item Explain the roles of \textbf{queries, keys, and values} in attention.
  \item Describe \textbf{multi-head attention} and why it's useful.
  \item Explain the components of a \textbf{transformer block} (attention, FFN, layer norm, residuals).
  \item Explain \textbf{causal masking} and why it's needed for language modelling.
  \item Describe why transformers need \textbf{positional encodings}.
  \item Explain the \textbf{quadratic complexity} of attention w.r.t.\ sequence length.
  \item Describe \textbf{weight tying} between embedding and unembedding layers.
  \item Explain \textbf{scaling laws}: how loss scales with parameters, data, and compute.
\end{itemize}

\bigseparator

% -----------------------------------------------------------------------------
%                 5.2 OVERVIEW + CORE CONCEPTS + EXTENDED THEORY
% -----------------------------------------------------------------------------
\subsection{Overview + Core Concepts + Extended Theory}

\subsubsection*{Roadmap}
\begin{enumerate}
  \item What are Large Language Models?
  \item Three LLM architectures: Decoder, Encoder, Encoder-Decoder
  \item Pretraining: self-supervision, cross-entropy loss, teacher forcing
  \item Training data and its problems
  \item Evaluating LLMs: perplexity
  \item Static vs contextual embeddings
  \item Self-attention mechanism
  \item Multi-head attention
  \item The transformer block
  \item Parallelising attention and causal masking
  \item Positional encodings
  \item Language modelling head
  \item Scaling laws and KV cache
\end{enumerate}

\bigseparator

% --- 5.2.1 Introduction to LLMs ---
\subsubsection{Introduction to Large Language Models}

\begin{defbox}[Large Language Model (LLM)]
A \textbf{Large Language Model} is a neural network that:
\begin{itemize}
  \item Takes a \textbf{context/prefix} as input
  \item Outputs a \textbf{probability distribution} over possible next words
  \item Has \textbf{billions of parameters}
  \item Is trained on \textbf{massive corpora} (web-scale data)
\end{itemize}

LLMs can \textbf{generate text} by repeatedly sampling from the next-word distribution.
\end{defbox}

\begin{keybox}[Continuity with Earlier LMs]
LLMs share the same fundamental task as n-gram and RNN language models:
\begin{itemize}
  \item Assign probabilities to sequences of words
  \item Generate text by predicting the next word
  \item Trained to minimise cross-entropy loss
\end{itemize}

\textbf{What's different:} Scale (billions of parameters), architecture (transformers), and training data (web-scale corpora).
\end{keybox}

\begin{tipbox}[Fundamental Intuition]
\textbf{Text contains enormous amounts of knowledge.}

Pretraining on lots of text with all that knowledge is what gives LLMs their remarkable abilities.
\end{tipbox}

\bigseparator

% --- 5.2.2 Three LLM Architectures ---
\subsubsection{Three LLM Architectures}

\begin{center}
\renewcommand{\arraystretch}{1.4}
\begin{tabular}{@{} l p{4cm} p{5.5cm} @{}}
\toprule
\textbf{Architecture} & \textbf{Examples} & \textbf{Characteristics} \\
\midrule
\textbf{Decoder-only} & GPT, Claude, Llama, Mistral, DeepSeek & Autoregressive (left-to-right); generates one token at a time; what people typically mean by ``LLM'' \\
\textbf{Encoder-only} & BERT, HuBERT & Masked language model (MLM); sees context on both sides; typically fine-tuned for classification \\
\textbf{Encoder-Decoder} & T5, Whisper, Flan-T5 & Maps input sequence to output sequence; used for translation, speech recognition, summarisation \\
\bottomrule
\end{tabular}
\end{center}

\begin{tipbox}[Decoder vs Encoder]
\begin{itemize}
  \item \textbf{Decoder:} Causal/autoregressive --- can only see past tokens (left-to-right generation).
  \item \textbf{Encoder:} Bidirectional --- can see tokens on both sides (used for understanding/classification).
\end{itemize}
\end{tipbox}

\bigseparator

% --- 5.2.3 Pretraining LLMs ---
\subsubsection{Pretraining Large Language Models}

\begin{defbox}[Three Stages of LLM Training]
\begin{enumerate}
  \item \textbf{Pretraining:} Learn to predict next word on massive unlabelled text.
  \item \textbf{Instruction Tuning:} Fine-tune on instruction-following examples.
  \item \textbf{Preference Alignment:} Align with human values and safety norms (e.g., RLHF).
\end{enumerate}

This week focuses on \textbf{pretraining}.
\end{defbox}

\begin{defbox}[Self-Supervised Training]
\textbf{Self-supervision:} The text itself provides the training signal.
\begin{enumerate}
  \item Take a corpus of text.
  \item At each position $t$, ask the model to predict the next word $w_{t+1}$.
  \item Train using gradient descent to minimise prediction error.
\end{enumerate}

\textbf{No manual labels required} --- the next word is the label!
\end{defbox}

\begin{defbox}[Cross-Entropy Loss for LM Training]
The loss at position $t$ is the negative log probability of the true next word:
\[
L_{CE}(\hat{y}_t, y_t) = -\log \hat{y}_t[w_{t+1}]
\]

\textbf{Intuition:}
\begin{itemize}
  \item If model assigns low probability to the correct word $\Rightarrow$ high loss.
  \item Training pushes the model to assign higher probability to correct words.
\end{itemize}

Since $y_t$ is one-hot, only the probability of the correct word matters.
\end{defbox}

\begin{defbox}[Teacher Forcing (Reminder)]
During training:
\begin{itemize}
  \item At each position $t$, compute loss for predicting $w_{t+1}$.
  \item Move to position $t+1$ using the \textbf{correct word} $w_{t+1}$, not the model's prediction.
\end{itemize}

\textbf{Why:} Prevents error accumulation during training.
\end{defbox}

\bigseparator

% --- 5.2.4 Training Data ---
\subsubsection{Training Data and Its Problems}

\textbf{Common pretraining corpora:}
\begin{itemize}
  \item \textbf{Common Crawl:} Snapshots of the entire web (billions of pages).
  \item \textbf{C4 (Colossal Clean Crawled Corpus):} 156 billion tokens, filtered English.
  \item \textbf{The Pile:} Curated mix of academic, web, books, and dialogue.
\end{itemize}

\begin{tipbox}[Problems with Web-Scraped Data]
\begin{itemize}
  \item \textbf{Copyright:} Much text is copyrighted; legality unclear.
  \item \textbf{Privacy:} Websites may contain private information (IP addresses, phone numbers).
  \item \textbf{Bias/Skew:} Disproportionately US-authored content; reflects societal biases.
  \item \textbf{Quality:} Boilerplate, adult content, duplicates need filtering.
  \item \textbf{Toxicity:} Toxicity detection can mistakenly flag minority dialects.
\end{itemize}
\end{tipbox}

\bigseparator

% --- 5.2.5 Evaluating LLMs: Perplexity ---
\subsubsection{Evaluating LLMs: Perplexity}

\begin{defbox}[Perplexity]
\textbf{Perplexity} measures how ``surprised'' the model is by unseen text.

For a test set of $n$ tokens $w_{1:n}$:
\[
\text{Perplexity}(w_{1:n}) = P(w_{1:n})^{-\frac{1}{n}} = \sqrt[n]{\prod_{i=1}^{n} \frac{1}{P(w_i \mid w_{<i})}}
\]

\textbf{Interpretation:}
\begin{itemize}
  \item \textbf{Lower perplexity = better model} (less surprised).
  \item Range: $[1, \infty)$.
  \item Minimising perplexity $\Leftrightarrow$ maximising probability.
\end{itemize}
\end{defbox}

\begin{tipbox}[Perplexity Comparisons]
Perplexity is sensitive to tokenisation, so \textbf{only compare models with the same tokeniser}.
\end{tipbox}

\textbf{Other evaluation factors:}
\begin{itemize}
  \item \textbf{Size:} Memory, training time, inference cost.
  \item \textbf{Energy:} kWh or CO$_2$ emissions.
  \item \textbf{Fairness:} Stereotypes, performance disparities across groups.
\end{itemize}

\bigseparator

% --- 5.2.6 Static vs Contextual Embeddings ---
\subsubsection{Static vs Contextual Embeddings}

\begin{tipbox}[Problem with Static Embeddings]
\textbf{Word2vec embeddings are static:} The same vector for a word regardless of context.

\textbf{Example:} ``The chicken didn't cross the road because \textbf{it} was too tired.''

What does ``it'' refer to? The static embedding can't distinguish:
\begin{itemize}
  \item ``it'' = chicken (tired)
  \item ``it'' = road (wide)
\end{itemize}
\end{tipbox}

\begin{defbox}[Contextual Embeddings]
\textbf{Contextual embedding:} A word's representation \textbf{changes depending on surrounding words}.

\textbf{Solution:} Use \textbf{attention} to build up the embedding by selectively integrating information from neighbouring words.

We say a word ``attends to'' some neighbours more than others.
\end{defbox}

\bigseparator

% --- 5.2.7 Self-Attention Mechanism ---
\subsubsection{Self-Attention Mechanism}

\begin{defbox}[Attention: Definition]
\textbf{Attention} is a mechanism for computing contextual embeddings by:
\begin{enumerate}
  \item Computing \textbf{similarity scores} between the current token and all context tokens.
  \item Using these scores to compute a \textbf{weighted sum} of context representations.
\end{enumerate}
\end{defbox}

\begin{defbox}[Query, Key, Value Roles]
Each token $x_i$ plays three roles, represented by learned projections:
\begin{align*}
q_i &= x_i W^Q & \text{(query: ``what am I looking for?'')} \\
k_i &= x_i W^K & \text{(key: ``what do I contain?'')} \\
v_i &= x_i W^V & \text{(value: ``what information do I provide?'')}
\end{align*}

\textbf{Matrices:} $W^Q, W^K \in \mathbb{R}^{d \times d_k}$, $W^V \in \mathbb{R}^{d \times d_v}$
\end{defbox}

\begin{thmbox}[Scaled Dot-Product Attention]
For current token $i$ with query $q_i$, attending to tokens $j \leq i$:

\textbf{1. Compute similarity scores:}
\[
\text{score}(x_i, x_j) = \frac{q_i \cdot k_j}{\sqrt{d_k}}
\]

\textbf{2. Normalise with softmax:}
\[
\alpha_{ij} = \text{softmax}(\text{score}(x_i, x_j)) = \frac{\exp(\text{score}(x_i, x_j))}{\sum_{k \leq i} \exp(\text{score}(x_i, x_k))}
\]

\textbf{3. Weighted sum of values:}
\[
a_i = \sum_{j \leq i} \alpha_{ij} v_j
\]
\end{thmbox}

\begin{tipbox}[Why Scale by $\sqrt{d_k}$?]
Dot products can become very large for high-dimensional vectors. Large values cause softmax to saturate (outputs near 0 or 1), leading to vanishing gradients. Scaling by $\sqrt{d_k}$ keeps values in a reasonable range.
\end{tipbox}

\bigseparator

% --- 5.2.8 Multi-Head Attention ---
\subsubsection{Multi-Head Attention}

\begin{defbox}[Multi-Head Attention]
Instead of one attention head, use $h$ heads in parallel:
\begin{itemize}
  \item Each head has its own $W^Q_c, W^K_c, W^V_c$ matrices.
  \item Each head can attend to different aspects of the context.
\end{itemize}

\textbf{Equations:}
\begin{align*}
\text{head}_c &= \text{Attention}(Q_c, K_c, V_c) \\
\text{MultiHead}(x_i) &= [\text{head}_1; \text{head}_2; \ldots; \text{head}_h] W^O
\end{align*}

where $W^O \in \mathbb{R}^{h \cdot d_v \times d}$ projects concatenated heads back to dimension $d$.
\end{defbox}

\begin{tipbox}[Why Multiple Heads?]
Different heads can learn to attend to different linguistic patterns:
\begin{itemize}
  \item One head might attend to syntactic dependencies.
  \item Another might attend to semantic relationships.
  \item Another might attend to nearby tokens.
\end{itemize}
\end{tipbox}

\bigseparator

% --- 5.2.9 The Transformer Block ---
\subsubsection{The Transformer Block}

\begin{defbox}[Transformer Block Components]
A transformer block contains:
\begin{enumerate}
  \item \textbf{Multi-Head Attention}
  \item \textbf{Residual Connection} (add input to output)
  \item \textbf{Layer Normalisation}
  \item \textbf{Feedforward Network} (2-layer MLP with ReLU)
  \item \textbf{Another Residual Connection + Layer Norm}
\end{enumerate}
\end{defbox}

\begin{defbox}[Transformer Block Equations]
\begin{align*}
t^1_i &= \text{LayerNorm}(x_i) \\
t^2_i &= \text{MultiHeadAttention}(t^1_i, [x_1, \ldots, x_N]) \\
t^3_i &= t^2_i + x_i & \text{(residual connection)} \\
t^4_i &= \text{LayerNorm}(t^3_i) \\
t^5_i &= \text{FFN}(t^4_i) \\
h_i &= t^5_i + t^3_i & \text{(residual connection)}
\end{align*}
\end{defbox}

\begin{defbox}[Layer Normalisation]
Normalise each vector to have zero mean and unit variance:
\[
\hat{x} = \frac{x - \mu}{\sigma}
\]
where $\mu = \frac{1}{d}\sum_i x_i$ and $\sigma = \sqrt{\frac{1}{d}\sum_i (x_i - \mu)^2}$

With learnable gain $g$ and bias $b$:
\[
\text{LayerNorm}(x) = g \cdot \frac{x - \mu}{\sigma} + b
\]
\end{defbox}

\begin{defbox}[Feedforward Network]
A 2-layer MLP applied independently to each position:
\[
\text{FFN}(x) = \text{ReLU}(x W_1 + b_1) W_2 + b_2
\]

Typically $d_{ff} = 4d$ (e.g., $d = 512$, $d_{ff} = 2048$).
\end{defbox}

\begin{tipbox}[Residual Stream View]
Attention is the \textbf{only} component that moves information between tokens. All other components operate on individual residual streams. Attention can be viewed as ``moving information from one token's stream to another.''
\end{tipbox}

\bigseparator

% --- 5.2.10 Parallelising Attention and Causal Masking ---
\subsubsection{Parallelising Attention and Causal Masking}

\begin{defbox}[Matrix Form of Attention]
Pack all $N$ tokens into matrix $X \in \mathbb{R}^{N \times d}$:
\begin{align*}
Q &= X W^Q, \quad K = X W^K, \quad V = X W^V \\
A &= \text{softmax}\left(\frac{Q K^\top}{\sqrt{d_k}}\right) V
\end{align*}

This computes attention for \textbf{all tokens simultaneously}!
\end{defbox}

\begin{defbox}[Causal Masking]
\textbf{Problem:} $QK^\top$ computes scores for \textbf{all} pairs, including future tokens.

\textbf{Solution:} Add a mask matrix $M$ where $M_{ij} = -\infty$ for $j > i$:
\[
A = \text{softmax}\left(\text{mask}\left(\frac{Q K^\top}{\sqrt{d_k}}\right)\right) V
\]

After softmax, positions with $-\infty$ become 0 (no attention to future).
\end{defbox}

\begin{tipbox}[Quadratic Complexity]
Attention is \textbf{$O(N^2)$} in sequence length:
\begin{itemize}
  \item $QK^\top$ is an $N \times N$ matrix.
  \item Every token attends to every other token.
\end{itemize}

This limits context length (though techniques like sparse attention help).
\end{tipbox}

\bigseparator

% --- 5.2.11 Positional Encodings ---
\subsubsection{Positional Encodings}

\begin{tipbox}[Why Transformers Need Position Information]
Attention treats all positions equally --- it has \textbf{no inherent notion of order}.

Without positional encodings, ``the cat sat on the mat'' and ``mat the on sat cat the'' would have the same representation!
\end{tipbox}

\begin{defbox}[Absolute Position Embeddings]
Learn a position embedding matrix $E_{\text{pos}} \in \mathbb{R}^{N_{\max} \times d}$:
\begin{itemize}
  \item One learned embedding for each position $1, 2, \ldots, N_{\max}$.
  \item Add position embedding to token embedding: $x_i = E[w_i] + E_{\text{pos}}[i]$
\end{itemize}
\end{defbox}

\bigseparator

% --- 5.2.12 Language Modelling Head ---
\subsubsection{Language Modelling Head}

\begin{defbox}[From Hidden State to Next-Word Probabilities]
After $L$ transformer layers, the final hidden state $h^L_N$ is transformed into a probability distribution:

\textbf{1. Unembedding layer:} Project to vocabulary-sized logits:
\[
u = h^L_N E^\top \quad \text{(shape: } 1 \times |V| \text{)}
\]

\textbf{2. Softmax:} Convert logits to probabilities:
\[
y = \text{softmax}(u)
\]

$y_k = P(w_{N+1} = k \mid w_{1:N})$
\end{defbox}

\begin{defbox}[Weight Tying]
The unembedding matrix is the \textbf{transpose of the embedding matrix} $E$:
\begin{itemize}
  \item \textbf{Embedding:} $E$ maps one-hot $\to$ embedding (shape $|V| \times d$)
  \item \textbf{Unembedding:} $E^\top$ maps embedding $\to$ logits (shape $d \times |V|$)
\end{itemize}

\textbf{Benefits:} Fewer parameters, often better generalisation.
\end{defbox}

\bigseparator

% --- 5.2.13 Scaling Laws and KV Cache ---
\subsubsection{Scaling Laws and KV Cache}

\begin{defbox}[Scaling Laws]
LLM performance (loss $L$) scales as a \textbf{power law} with:
\begin{itemize}
  \item \textbf{Model size $N$} (number of parameters)
  \item \textbf{Dataset size $D$} (tokens)
  \item \textbf{Compute $C$} (FLOPs)
\end{itemize}

\[
L(N) = \left(\frac{N_c}{N}\right)^{\alpha_N}, \quad L(D) = \left(\frac{D_c}{D}\right)^{\alpha_D}, \quad L(C) = \left(\frac{C_c}{C}\right)^{\alpha_C}
\]

\textbf{Key insight:} We can predict final performance from early training!
\end{defbox}

\begin{exbox}[Estimating Parameters]
For a transformer with $n$ layers and dimension $d$:
\[
N \approx 12 \cdot n \cdot d^2
\]

GPT-3: $n = 96$ layers, $d = 12288$ $\Rightarrow$ $N \approx 175$ billion parameters.
\end{exbox}

\begin{defbox}[KV Cache]
\textbf{Problem:} At inference, we generate tokens one at a time. Recomputing all keys and values at each step is wasteful.

\textbf{Solution:} Cache the key and value vectors for all previous tokens.

At step $i$:
\begin{itemize}
  \item Compute $q_i, k_i, v_i$ for the new token.
  \item Retrieve $k_1, \ldots, k_{i-1}$ and $v_1, \ldots, v_{i-1}$ from cache.
  \item Compute attention using cached keys/values.
  \item Store $k_i, v_i$ in cache for future steps.
\end{itemize}
\end{defbox}

\bigseparator

% --- Summary Table ---
\subsubsection{Summary: Transformers vs RNNs}

\begin{center}
\renewcommand{\arraystretch}{1.3}
\begin{tabular}{@{} l p{5cm} p{5cm} @{}}
\toprule
& \textbf{RNN/LSTM} & \textbf{Transformer} \\
\midrule
Processing & Sequential (one step at a time) & Parallel (all tokens at once) \\
Context & Via hidden state (limited) & Via attention (direct access) \\
Long-range deps & Difficult (vanishing gradients) & Easy (direct attention) \\
Training & Sequential backprop through time & Highly parallelisable \\
Complexity & $O(N)$ per step & $O(N^2)$ for attention \\
Position info & Implicit (order of processing) & Explicit (positional encodings) \\
\bottomrule
\end{tabular}
\end{center}

